\chapter{Podstawy metodyczne}

% Odpowiednik sekcji Methodological Background z czasopism poświęconych
% oprogramowaniu badawczemu. Opisujesz matematykę i logikę, które
% zaimplementowałeś – jeszcze nie kod. Objętość: 8-12 stron.
%
% Zasada nadrzędna tego rozdziału: po każdym wzorze następuje akapit
% wyjaśniający, co ten wzór robi i po co, językiem zrozumiałym dla badacza
% bez przygotowania matematycznego. Wzór bez wyjaśnienia jest w tej pracy
% błędem merytorycznym, nie oszczędnością miejsca.

Wprowadź dziedzinę i umieść w niej metodę. Wyjaśnij, jaki problem decyzyjny, pomiarowy albo przyczynowy metoda rozwiązuje i czym różni się od podejść wcześniejszych. Odwołaj się do artykułu źródłowego oraz do dwóch, trzech opracowań wprowadzających.

\section{Aparat formalny}

Przedstaw wielkości wejściowe i ich dziedziny, a następnie kolejne kroki algorytmu. Wzory, do których praca się odwołuje, numeruj i etykietuj:

\begin{equation}\label{eq:przyklad}
  w_j = \frac{d_j}{\sum_{k=1}^{n} d_k}, \qquad d_j = 1 - e_j.
\end{equation}

We wzorze \eqref{eq:przyklad} symbol $w_j$ oznacza wagę kryterium $j$, $e_j$ – jego entropię, a $n$ – liczbę kryteriów. Objaśnij w ten sposób każdy symbol, którego używasz.

\section{Założenia i granice stosowalności}

Wypisz założenia, przy których metoda jest poprawna: rozkłady, niezależność obserwacji, kompletność danych, skale pomiarowe. Następnie napisz, co się dzieje, gdy założenie jest naruszone. Ta sekcja zasila później kontrakt w \plik{SPEC.md} oraz testy odporności pakietu.

\section{Rozwiązania pokrewne}

Przegląd istniejących implementacji w układzie problemowym: co już jest dostępne, czego brakuje, jak Twoje narzędzie się do tego ma. Zakończ jednoznacznym stwierdzeniem, jaki wkład wnosi Twój pakiet.
